%% ma: L10-B14-0016
%% lop: 10
%% chuong: 5
%% bai: 14
%% dang: 10.14.1
%% loai: TN4
%% muc: NB
%% dap_an: "D"
%% nguon: "chuyen-de-cac-so-dac-trung-cua-mau-so-lieu-khong-ghep-nhom-toan-10-chuong-trinh-moi.pdf (GV. Phan Nhật Linh) (P04), Bài 03 Dạng 1 Phần I Câu 6"
%% kien_thuc: "Lớp 10 Bài 14 (khoảng biến thiên)"
%% trang_thai: chinh-thuc
%% ghi_chu: "Tự duyệt theo yêu cầu giáo viên 22/10/2026."
\begin{ex}
Số sản phẩm sản xuất mỗi ngày của một phân xưởng trong $9$ ngày liên tiếp được ghi lại như sau:
\begin{center}
\begin{tabular}{|c|c|c|c|c|c|c|c|c|}\hline
$27$ & $26$ & $21$ & $28$ & $25$ & $30$ & $26$ & $23$ & $26$\\\hline
\end{tabular}
\end{center}
Khoảng biến thiên của mẫu số liệu là
\choice{$8$}{$5$}{$6$}{\True $9$}
\loigiai{$R=30-21=9$. Chọn D.}
\end{ex}
